\pdfbookmark[1]{II. Bestimmung der nächsten ganzen Zahl}{complexe-II}
\begin{center}[II.]\par\medskip
BESTIMMUNG DER NÄCHSTEN GANZEN ZAHL.
\end{center}
Es sei
\[
\begin{aligned}
\varepsilon^3&=1,\quad m=a+b\varepsilon+c\varepsilon\varepsilon\\
2a-b-c&=A+\alpha\\
2b-c-a&=B+\beta\\
2c-a-b&=C+\gamma
\end{aligned}
\]
wobei $A,B,C$ ganze Zahlen und $\alpha,\beta,\gamma$ positive echte Brüche sind. Dann gilt
\[
A+B+C+\alpha+\beta+\gamma=0
\]
\Needspace{5\baselineskip}
Es sind also drei Fälle zu unterscheiden:

I. $\alpha+\beta+\gamma=0$, folglich $\alpha=0$, $\beta=0$, $\gamma=0$

1, $A\equiv B\equiv C\ (\text{mod. }3)$. Hier ist $m$ selbst eine ganze Zahl.

2, $A-B\equiv B-C\equiv C-A\equiv\pm1\ (\text{mod. }3)$. Hier ist $m\pm\frac{\varepsilon-\varepsilon\varepsilon}{3}\pdot\varepsilon^n$ eine ganze Zahl.

II. $\alpha+\beta+\gamma=1$.

Hier ist jede der Zahlen
\[
\begin{aligned}
&A+B\varepsilon+C\varepsilon\varepsilon+1\\
&A+B\varepsilon+C\varepsilon\varepsilon+\varepsilon\\
&A+B\varepsilon+C\varepsilon\varepsilon+\varepsilon\varepsilon
\end{aligned}
\]
\SourceLeaf{402}%
durch $1-\varepsilon$ teilbar und eine davon durch $3$. Der Quotient oder
\[
m+\frac{\varepsilon^n-\alpha-\beta\varepsilon-\gamma\varepsilon\varepsilon}{3}
\]
ist die gesuchte ganze Zahl.

III. $\alpha+\beta+\gamma=2$

Hier sind
\[
\begin{aligned}
&A+B\varepsilon+C\varepsilon\varepsilon+\varepsilon+\varepsilon\varepsilon\\
&A+B\varepsilon+C\varepsilon\varepsilon+\varepsilon\varepsilon+1\\
&A+B\varepsilon+C\varepsilon\varepsilon+1+\varepsilon
\end{aligned}
\]
durch $1-\varepsilon$ und eine dieser Zahlen durch $3$ teilbar. Der Quotient oder
\[
m+\frac{\varepsilon^n(\varepsilon+\varepsilon\varepsilon)-\alpha-\beta\varepsilon-\gamma\varepsilon\varepsilon}{3}
\]
ist die gesuchte ganze Zahl.

In allen drei Fällen hat der Rest die Form
\[
x+y\varepsilon+z\varepsilon\varepsilon
\]
sodass $x,y,z$ dem Betrag nach kleiner als $\frac13$ sind und $x+y+z=0$ gilt. Daraus folgt notwendig
\[
xx+yy+zz=2xx-2yz=2yy-2xz=2zz-2xy<\tfrac29
\]
weil zwei der drei Größen $x,y,z$ notwendig dasselbe Vorzeichen haben. Folglich ist die Determinante des Restes
\[
=\tfrac32(xx+yy+zz)<\tfrac13\qquad\text{Q. E. D.}
\]

Die nächste ganze Zahl lässt sich \emph{bequemer} auf folgende Weise bestimmen. Gegeben sei $a+b\varepsilon+c\varepsilon\varepsilon=m$; man setze
\[
\begin{aligned}
b-a&=C+\gamma\\
c-b&=A+\alpha\\
a-c&=B+\beta
\end{aligned}
\]
wobei $A,B,C$ die nächstkleineren ganzen Zahlen und $\alpha,\beta,\gamma$ positive Brüche sind.
Hier sind drei Fälle zu unterscheiden:

\SourceLeaf{403}%
I. $\alpha+\beta+\gamma=0$; dann ist $m$ selbst eine ganze Zahl

II. $\alpha+\beta+\gamma=1$; dann ist die nächste ganze Zahl
\[
\begin{array}{ll}
B+(B+C)\varepsilon&\text{wenn }\alpha\text{ der größte Bruch ist.}\\
\qquad C\varepsilon+(A+C)\varepsilon\varepsilon&\qquad\beta\\
A+B\quad\pdot\quad+A\varepsilon\varepsilon&\qquad\gamma
\end{array}
\]

II. $\alpha+\beta+\gamma=2$; dann ist die nächste ganze Zahl
\[
\begin{array}{ll}
B+1+(B+C+2)\varepsilon&\text{wenn }\alpha\text{ der kleinste Bruch ist.}\\
\qquad(C+1)\varepsilon+(A+C+2)\varepsilon\varepsilon&\qquad\beta\\
A+B+2\quad\pdot\quad+(A+1)\varepsilon\varepsilon&\qquad\gamma
\end{array}
\]
In II, $1$ ist der Rest $\beta+(\beta+\gamma)\varepsilon$, dessen Determinante
\[
\begin{aligned}
&=\beta\beta+\beta\gamma+\gamma\gamma\\
&=\tfrac13-\tfrac16[(\alpha-\beta)(1+3\beta)+(\alpha-\gamma)(1+3\gamma)]
\end{aligned}
\]
ist.

\emph{Noch einfacher so:}

Man ordne die Brüche $a-[a],b-[b],c-[c]$ nach ihrer Größe und bezeichne sie der Reihe nach mit $p,q,r$. Sind alle drei gleich groß, so ist $m$ eine ganze Zahl. Sind sie aber ungleich, so sei $t$ ein beliebiger Bruch zwischen
\[
\begin{array}{llll}
p\text{ und }q,&\text{jenachdem}&q-p&\text{am größten ist}\\
q\text{ und }r&&r-q&\\
r\text{ und }1+p&&1+p-r&
\end{array}
\]
Dann ist
\[
[a-t]+[b-t]\varepsilon+[c-t]\varepsilon\varepsilon
\]
die nächste ganze Zahl.
\BookRule{.3\textwidth}

\SourceLeaf{404}%
\pdfbookmark[1]{III.}{complexe-III}
\begin{center}[III.]\end{center}
Es sei $\varepsilon^5=1$
\[
\begin{aligned}
a+b\varepsilon+c\varepsilon\varepsilon+d\varepsilon^2+e\varepsilon^4&=q'\\
a+b\varepsilon^{-1}+c\varepsilon^{-2}+d\varepsilon^{-3}+e\varepsilon^{-4}&=q''''
\end{aligned}
\]
\[
\begin{aligned}
(a-b)^2+(b-c)^2+(c-d)^2+(d-e)^2+(e-a)^2&=2p'\\
(a-c)^2+(b-d)^2+(c-e)^2+(d-a)^2+(e-b)^2&=2p''
\end{aligned}
\]
\[
\begin{aligned}
q'q''''&=-p'\varepsilon-p''\varepsilon\varepsilon-p''\varepsilon^3-p'\varepsilon^4=P'\\
q''q'''&=-p'\varepsilon\varepsilon-p''\varepsilon^4-p''\varepsilon-p'\varepsilon^3=P''
\end{aligned}
\]

\[
\text{Determinante}=P'P''=-p'p'+3p'p''-p''p''
\]
\[
\begin{aligned}
\text{Mensura}&=2p'+2p''=2P'+2P''\\
&=5(aa+bb+cc+dd+ee)-(a+b+c+d+e)^2
\end{aligned}
\]
Multiplikation mit $1-\varepsilon$ ergibt die neue Mensur $=8p'$
\[
\text{Höchste Mensur}=2\left(\frac{\sin72^\circ}{\sin36^\circ}+\frac{\sin36^\circ}{\sin72^\circ}\right)\surd D=4{,}472\surd D
\]

\emph{Modul} $=1-\varepsilon$
\[
\begin{aligned}
1-\varepsilon&=x\\
\varepsilon&=1-x\\
\varepsilon\varepsilon&=1-2x+xx\\
\varepsilon^3&=1-3x+3xx-x^3\\
\varepsilon^4&=1-4x+6xx-4x^3+x^4\\
&=-4+6x-4xx+x^3
\end{aligned}
\]
Also
\[
\begin{aligned}
\frac{1-\varepsilon^n}{1-\varepsilon}&\equiv n\text{ mod. }(1-\varepsilon)\\
\varepsilon^n&\equiv1-nx\text{ mod. }(1-\varepsilon)^2\\
\left(\frac{\varepsilon+\varepsilon^4}{\varepsilon\varepsilon+\varepsilon^3}\right)^n&\equiv1+nxx\text{ mod. }(1-\varepsilon)^3
\end{aligned}
\]
Eine Zahl, die $\equiv1\text{ mod. }(1-\varepsilon)^3$ ist, kann also \emph{nur} dann eine Einheit sein, wenn sie zugleich $\equiv1\ (\text{mod. }5)$ ist.
\BookRule{.3\textwidth}

\SourceLeaf{405}%
\Needspace{17\baselineskip}
\pdfbookmark[1]{IV. Einiges über die Mensur der Zahlen}{complexe-IV}
\begin{center}[IV.]\par\medskip
EINIGES ÜBER DIE MENSUR DER ZAHLEN.
\end{center}
Es sei $\varepsilon^n=1$, wobei $n$ eine Primzahl ist,
\[
\begin{aligned}
m&=a+a'\varepsilon+a''\varepsilon\varepsilon+a'''\varepsilon^3+\ldots+a^{(n-1)}\varepsilon^{n-1}=f\varepsilon\\
D&=f\varepsilon\pdot f\varepsilon\varepsilon\pdot f\varepsilon^3\ldots f\varepsilon^{n-1}\\
f\varepsilon\pdot f\varepsilon^{n-1}&=-b'(\varepsilon+\varepsilon^{n-1})-b''(\varepsilon^2+\varepsilon^{-2})\\
&\qquad-b'''(\varepsilon^3+\varepsilon^{-3})\ldots
\end{aligned}
\]
dann ist
\[
\begin{aligned}
2b'&=(a-a')^2+(a'-a'')^2+(a''-a''')^2+\text{usw.}\\
2b''&=(a-a'')^2+(a'-a''')^2+(a''-a'''')^2+\text{usw.}\\
&\text{usw.}
\end{aligned}
\]

Hier sind also $b',b'',b'''\dotstwo$ ausschließlich positive Größen; sie heißen \emph{Partialmensuren} von $m$, ihre Summe
\[
\begin{aligned}
b'+b''+b'''+\text{usw.}={}&n(aa+a'a'+a''a''+\dotstwo)\\
&\qquad-(a+a'+a''+\text{usw.})^2
\end{aligned}
\]
die \emph{Generalmensur}. Setzt man
\[
f\varepsilon\pdot f\varepsilon^{n-1}=c',\quad f\varepsilon\varepsilon\pdot f\varepsilon^{n-2}=c''\quad\text{usw.}
\]
so ist
\[
\begin{aligned}
&c'+c''+c'''+\text{usw.}+c^{\frac12(n-1)}\\
&\qquad=b'+b''+b'''+\text{usw.}+b^{\frac12(n-1)}
\end{aligned}
\]
\[
\begin{aligned}
&c'(\varepsilon+\varepsilon^{n-1})+c''(\varepsilon\varepsilon+\varepsilon^{n-2})+c'''(\varepsilon^3+\varepsilon^{n-3})+\text{usw.}\\
&\qquad=2(b'+b''+b'''+\text{usw.}+b^{\frac12(n-1)})-nb'
\end{aligned}
\]
\[
\begin{aligned}
&c'(2-\varepsilon-\varepsilon^{n-1})+c''(2-\varepsilon\varepsilon-\varepsilon^{n-2})\\
&\qquad+c'''(2-\varepsilon^3-\varepsilon^{n-3})+\text{usw.}=nb'
\end{aligned}
\]
\[
b'>\frac{n-1}{2n}(nD)^{\frac2{n-1}},\quad b'+b''+b'''+\text{usw.}>\frac{n-1}{2}\pdot D^{\frac2{n-1}}
\]

Gilt allgemein
\[
f\varepsilon\pdot f\varepsilon^{n-1}=A+A'\varepsilon+A''\varepsilon\varepsilon+A'''\varepsilon^3+\dotstwo
\]
so ist die Generalmensur $\triangle=-A-A'-A''-\text{usw.}+nA$

Mensur von $(1+\varepsilon)f\varepsilon\ldots$
\[
\triangle'=4\triangle-2n(A-A')=4\triangle-2nb'
\]
Ist $a+a'+a''+\dotstwo=0$, so ist $\triangle=n(aa+a'a'+a''a''+\text{usw.})$

und ist $A+A'+A''+\ldots=0$, so ist $\triangle=nA$, $\triangle'=n(2A+2A')$

Ist also einer der Koeffizienten $A',A''$ usw. negativ und dem Betrag nach größer als $\frac12A$, so lässt sich die Mensur verringern, ohne die Determinante zu ändern.
\BookRule{.3\textwidth}

\SourceLeaf{406}%
\Needspace{13\baselineskip}
\pdfbookmark[1]{V.}{complexe-V}
\begin{center}[V.]\end{center}
Sollte sich bestätigen, dass jede Einheit nur aus Faktoren der Form
\[
\frac{\varepsilon^\alpha-\varepsilon^\beta}{\varepsilon^\gamma-\varepsilon^\delta}
\]
zusammengesetzt ist, so wäre folgender Satz bewiesen:

\emph{Ist $f(\varepsilon)$ eine Einheit, so ist}
\[
\frac{f(\varepsilon)}{f(\varepsilon^{-1})}=\varepsilon^n
\]

Auch ohne \emph{jenen} Satz \emph{vorauszusetzen}, lässt sich der Schlusssatz leicht beweisen.
Es sei
\[
\frac{f\varepsilon}{f\varepsilon^{-1}}=F\varepsilon
\]
dann ist
\[
F\varepsilon\pdot F\varepsilon^{-1}=1
\]
woraus sich mithilfe der Lehre von der Mensur leicht ergibt, dass
\[
F\varepsilon=\pm\varepsilon^n
\]
Das untere Zeichen ist jedoch unmöglich, weil sonst $f\varepsilon$ durch $1-\varepsilon$ teilbar sein müsste.

Dass die Determinante einer von $0$ verschiedenen Zahl nicht $=0$ sein kann, lässt sich leicht beweisen. Ist die Determinante durch $m$ teilbar, so ist die Zahl selbst durch $1-\varepsilon$ teilbar. Ist also die Determinante durch $m^{m-1}$ teilbar, so muss die Zahl selbst durch $m$ teilbar sein. Dies ist widersinnig, denn bei der Determinante $0$ könnte man die Zahl zunächst so oft durch $m$ teilen, ohne die Determinante zu ändern, bis sie nicht mehr teilbar wäre. Der erste Satz wird so ersichtlich. Es sei die gegebene Zahl
\[
a+b\varepsilon+c\varepsilon\varepsilon+\text{usw.}\equiv a+b+c\dotstwo\text{ mod. }1-\varepsilon
\]
also ihre Determinante $\equiv(a+b+c\dotstwo)^{m-1}\text{ mod. }1-\varepsilon$.
\BookRule{.3\textwidth}

\SourceLeaf{407}%
\Needspace{16\baselineskip}
\pdfbookmark[1]{VI.}{complexe-VI}
\begin{center}[VI.]\end{center}
Es sei $\varepsilon^n=1$
\[
\begin{aligned}
f\varepsilon&=a+b\varepsilon+c\varepsilon\varepsilon+d\varepsilon^3+\text{usw.}\\
m&=\text{Determinante dieser Zahl}\\
\frac{m}{f\varepsilon}&=f\varepsilon\varepsilon\pdot f\varepsilon^3\ldots f\varepsilon^{n-1}\\
&=A+B\varepsilon+C\varepsilon\varepsilon+\text{usw.}=F\varepsilon
\end{aligned}
\]
Der Zahl $f\varepsilon$ entspricht eine Wurzel der Kongruenz $x^n\equiv1\ (\text{mod. }m.)$ Sie sei $r$. Es gilt
\[
\begin{aligned}
nA&=F1+F\varepsilon+F\varepsilon\varepsilon+\dotstwo\\
nB&=F1+\varepsilon^{-1}F\varepsilon+\varepsilon^{-2}F\varepsilon\varepsilon+\dotstwo\\
nC&=F1+\varepsilon^{-2}F\varepsilon+\varepsilon^{-4}F\varepsilon\varepsilon+\dotstwo\\
&\text{usw.}
\end{aligned}
\]

Da $F\varepsilon\varepsilon,F\varepsilon^3,F\varepsilon^4$ usw. durch $f\varepsilon$ teilbar sind, sind also
\[
\begin{gathered}
nA-F1-\varepsilon(nB-F1)\\
nA-F1-\varepsilon\varepsilon(nC-F1)\\
nA-F1-\varepsilon^3(nD-F1)\\
\text{usw.}
\end{gathered}
\]
alle durch $f\varepsilon$ teilbar oder auch
\[
\begin{gathered}
n(A-B)-\varepsilon n(B-C)\\
n(B-C)-\varepsilon n(C-D)\\
\text{usw.}
\end{gathered}
\]
durch $f\varepsilon$ teilbar; folglich [wenn $f\varepsilon$ durch $1-\varepsilon$ und $F\varepsilon$ durch eine reelle ganze Zahl nicht teilbar ist]
\[
\varepsilon\equiv\frac{A-B}{B-C}\equiv\frac{B-C}{C-D}\equiv\frac{C-D}{D-E}\text{ usw. }(\text{mod. }f\varepsilon)
\]
\BookRule{.3\textwidth}

\clearpage
\SourceLeaf{408}%
\pdfbookmark[1]{Bemerkungen zu den komplexen Zahlen}{complexe-Bemerkungen}
\begin{center}BEMERKUNGEN.\end{center}
\BookRule{.15\textwidth}
Die hier unter der gemeinsamen Überschrift Zur Theorie der komplexen Zahlen zusammengestellten Untersuchungen sind verstreute Notizen in der Handschrift. Sie enthalten die wesentlichen Schritte des Beweises des \textsc{Fermat}schen Satzes für die dritte und fünfte Potenz. Die aus dritten Einheitswurzeln zusammengesetzten Zahlen sind in unvollständigen, hier nicht abgedruckten Aufzeichnungen sowohl mithilfe der Theorie der binären quadratischen Formen als auch der Kreisteilung untersucht. Bei der Anwendung der letzteren, und zwar während der Ausarbeitung der Abhandlung \emph{Disquisitionum circa aequationes puras ulterior evolutio}, wurde auch die ternäre kubische Form aufgestellt, in die $27\frac{x^n-1}{x-1}$ für eine Primzahl $n\equiv1\text{ mod. }3$ umgewandelt werden kann. Zugleich wurde die Theorie der Komposition der damit verwandten Form $X^3+mY^3+mmZ^3-3mXYZ$ entwickelt.

Die in den Untersuchungen des Bruchstücks [I] vorausgesetzte Eigenschaft der aus dritten Einheitswurzeln gebildeten ganzen Zahlen, dass jede nur auf eine Weise in Primfaktoren zerlegt werden kann, ergibt sich aus dem \textsc{Euklid}ischen Verfahren zur Bestimmung der gemeinsamen Teiler zweier Zahlen, wenn man dabei den unter [II] hergeleiteten Satz über die nächste ganze Zahl zu einem beliebig vorgegebenen Bruch anwendet.

Dass dieselbe grundlegende Eigenschaft auch den aus fünften Einheitswurzeln zusammengesetzten Zahlen zukommt, folgt daraus, dass der nach einer ganz entsprechenden Regel wie in [II] gebildete Bruchrest entweder von $m$ oder aber von $m$ multipliziert mit einer geeigneten Einheit $E$ so beschaffen ist, dass er durch Subtraktion von der gegebenen Zahl $mE$ eine ganze Zahl ergibt und dass seine Determinante die Einheit nicht übersteigt. Die Einheiten lassen sich aber, wie in [III] angedeutet, aus der Theorie der binären quadratischen Formen der Determinante $5$ in Verbindung mit der Zerlegung irgendeiner reellen Primzahl in vier Faktoren (zum Beispiel $11=\text{Det. }(2+\varepsilon)$) ableiten, nämlich als Produkte der Potenzen von $\varepsilon$ und $1+\varepsilon$.
\begin{flushright}\textsc{Schering.}\end{flushright}
\BookRule{.3\textwidth}
